\section{Discussion}
\label{sec:discussion}

The central claim of this work is that municipal PAM labels are enough to train a pixel-level productivity model whose intramunicipal maps are the intended product, while municipal means of those maps remain evaluable against official statistics.
On the 2021 validation and test splits, those aggregates attain roughly $0.35$--$0.36$\,\si{t/ha} RMSE and about $8\%$--$9\%$ MAPE, with closely matched $R^2$ near $0.23$.
MAPE suggests that relative municipal errors are often practically small, whereas $R^2$ indicates that only a modest fraction of cross-municipality variance is explained on this single holdout year.
Pixel maps (e.g., Figure~\ref{fig:guara}) cannot be scored directly without field labels; their credibility rests on consistency of the municipal mean and on qualitative spatial structure.

Adding Xavier climate covariates improves both validation and test $R^2$ relative to the spectral-only ablation (validation $R^2$ $0.227$ vs.\ $0.165$; test $R^2$ $0.228$ vs.\ $0.117$).
Climate inputs are therefore useful for the municipal selection criterion and for spatial generalization within the holdout year under the protocol used here.

Architecture search on spectral inputs favored a compact Transformer ($d_{\mathrm{model}}=96$, modest feed-forward width, moderate dropout) for the spectral-only ablation, whereas the proposed model uses a wider configuration ($d_{\mathrm{model}}=128$, $d_{\mathrm{inner}}=128$, dropout $0.3$).
Combined with validation degradation after the learning-rate schedule step, this points to a regime where capacity, covariates, and optimization noise jointly affect year-holdout generalization of the downscaled aggregates.

Per-year diagnostics highlight 2022 as a severe failure mode: RMSE and MAPE rise sharply even though $R^2$ is near zero.
Interpreting $R^2$ alone can therefore be misleading under drought-driven level shifts; RMSE and MAPE remain essential, and models selected on a single holdout year may not transfer to extreme seasons---including the intramunicipal maps derived from the same checkpoint.

NDVI-weighted temporal pooling improved validation $R^2$ relative to mean pooling over valid timesteps on the spectral-only architecture, but test $R^2$ was essentially unchanged.
The ablation supports NDVI weighting as a useful inductive bias for model selection, while cautioning against overinterpreting a single holdout-year test difference.

Early-season results show progressive information gain as the season unfolds.
Pooled across years, municipal $R^2$ turns positive by month 3 and peaks near month 5 ($0.429$) before a slight decline at month 6.
On 2021 alone, truncated evaluations remain weak until the full six-month window recovers the holdout $R^2$.
Operational early warning with downscaled maps would therefore need models trained explicitly for truncated inputs, multi-year aggregation of forecasts, or stronger season-aware regularization.

A structural limitation is the mean-aggregation assumption: municipal productivity is treated as the average of pixel contributions, which ignores possible unequal area weighting, management heterogeneity, and label noise in PAM.
Without pixel-level references, we cannot quantify intramunicipal calibration error, only municipal consistency.
Further limitations include a single holdout year, the absence of external algorithmic baselines, and sensitivity to extreme seasons such as 2022.
Self-supervised pretraining was not evaluated here and remains future work.
